\documentclass[11pt]{article}
\usepackage[margin=1in]{geometry}
\usepackage{amsmath,graphicx}
\usepackage[hidelinks]{hyperref}
\title{Running: Interval Training}
\author{Liam McGrath}
\date{4 October 2026}
\begin{document}
\maketitle

\section{Overview}
Interval training is intended to improve overall running pace and VO$_2$ max. Liam McGrath adopted this training method in 2013 during his marathon training, as part of his cross-country coaches' training plan, since they were training for a marathon themselves.

\section{Best times, 2013 and 2017}
Earlier history is missing, but these are rough estimates of the best times from the major training blocks of 2013 and 2017. Each time is the best single repeat; pace per kilometre is worked out from the repeat distance.

\begin{table}[htbp]
\centering
\begin{tabular}{llll}
\hline
Year & Session & Best repeat & Pace per km \\
\hline
2013 & $12 \times 400$ m & 1:08 & 2:50 \\
2017 & $8 \times 800$ m & 2:42 & 3:23 \\
2017 & $4 \times 1600$ m & 5:38 & 3:31 \\
\hline
\end{tabular}
\caption{Rough estimates of best repeat times, recalled rather than logged.}
\end{table}

\section{Routes}
Alternate routes were run between 2013 and 2026. One of them is I1, shown below: a 0.95 km loop around neighbourhood streets rather than a track. Times from those years have been lost to old cell phones.

\begin{figure}[htbp]
\centering
\includegraphics[width=0.9\linewidth]{tex_images/athletics/interval_track_1.png}
\caption{Interval route 1 (I1), measured at 0.95 km with GeoDistance. The start marker and map coordinates have been removed.}
\end{figure}

\section{The current log}
Sessions on I1 have been logged interval by interval since 22 September 2026. The Running page of the website charts every interval, each session's average and the pace per kilometre, and offers the log as JSON and CSV downloads.

\end{document}
